
% SGA 3，第 VII 篇_B；英文重建本第一轮的中文译本。
% 范围： 节 4.2.2 全文 依据文献局部页 75--76,
% 止于下节之前： 节 4.3.
% 合并阅读本页： 588--589. 重版印刷页： 577--578.
% 范围内原版边码： 544.
% 范围内编者注： 135--136.
% 依据： Polo--Gille Exp7B-13oct24.pdf,
% SHA-256 D303690E0195654BBE7C3DAA5D89EA23E24CD709A3575217234AC19AE8836694.
% 参照了 Jacob Reinhold 的 CC BY 4.0 Markdown，作为对照来源。
% 原英文 Loop 1 阶段说明：泛性质图为第一轮直接从 PDF 截取的原图。（当前中文工程已采用原生图。）
\subsubsection*{4.2.2}
\phantomsection
\label{sga3:VIIB:4.2.2}
设 \(\mathbf L\) 是 \(\mathbf O_k\) 上的平坦 \(p\)-Lie 代数。结合 VII\textsubscript{A}，\SGARefLink{sga3:VIIA:5.3.3}{5.3.3}，\SGARefLink{sga3:VIIB:2.6.2:prop}{命题 2.6.2} 第\textup{\SGARefLink{sga3:vii:paragraph:5-1:equation:i}{(i)}} 项的论证表明，\(U_p(\mathbf L)\) 在 \(\mathbf O_k\) 上平坦。因此，\(\mathscr G_p(\mathbf L)\) 是在 \(k\) 上拓扑平坦的形式 \(k\)-群，其协变 \(\mathbf O_k\)-双代数为 \(U_p(\mathbf L)\)。
\footnote{\label{sga3:VIIB:nde:135}以下细节由编者补充。}
由 \SGARefLink{sga3:VIIB:2.6.2:prop}{2.6.2}\textup{\SGARefLink{sga3:VIIB:2.6.2:prop:item:iii}{(iii)}} 的证明，对每个有限长度的 \(k\)-代数 \(C\)，Lie 代数 \(\operatorname{Lie}(\mathscr G_p(\mathbf L))(C)\) 是
\[
  U_p(\mathbf L)(C)
  =
  U_p\bigl(\mathbf L(C)\bigr);
\]
的本原元素集合；另见 VII\textsubscript{A}，\SGARefLink{sga3:VIIA:3.2.3}{3.2.3}。此外，由 VII\textsubscript{A}，\SGARefLink{sga3:VIIA:5.5.3}{5.5.3}，典范态射 \(\mathbf L\to U_p(\mathbf L)\) 诱导 \(p\)-Lie 代数同构
\[
  \tau_{p,\mathbf L}:
  \mathbf L
  \xrightarrow{\ \sim\ }
  \operatorname{Lie}\bigl(\mathscr G_p(\mathbf L)\bigr).
\]
可与特征 \(0\) 时的 \SGARefLink{sga3:VIIB:3.1:lemma}{3.1} 比较。
形式群 \(\mathscr G_p(\mathbf L)\) 可由一个泛性质刻画。每个到形式群 \(G\) 的态射
\[
  h:\mathscr G_p(\mathbf L)\longrightarrow G
\]
诱导一个态射
\[
  \operatorname{Lie}(h):
  \operatorname{Lie}\bigl(\mathscr G_p(\mathbf L)\bigr)
  \longrightarrow
  \operatorname{Lie}(G).
\]
将这个态射与 \(\tau_{p,\mathbf L}\) 复合，得到
\[
  h':\mathbf L\longrightarrow\operatorname{Lie}(G).
\]
\medskip
\noindent\textbf{命题。}
\phantomsection
\label{sga3:VIIB:4.2.2:universal-property}
{\itshape 设 \(k\) 是特征 \(p>0\) 的伪紧环，\(G\) 是形式 \(k\)-群，\(\mathbf L\) 是 \(\mathbf O_k\) 上的平坦 \(p\)-Lie 代数。上面定义的映射 \(h\mapsto h'\) 是双射
\[
  \operatorname{Hom}_{\operatorname{Grf}/k}
    \bigl(\mathscr G_p(\mathbf L),G\bigr)
  \xrightarrow{\ \sim\ }
  \operatorname{Hom}_{p\text{-}\operatorname{Lie}}
    \bigl(\mathbf L,\operatorname{Lie}(G)\bigr),
\]
其中右边是从 \(\mathbf L\) 到 \(\operatorname{Lie}(G)\) 的 \(p\)-Lie 代数态射的集合。}
\textit{证明。}\footnote{\label{sga3:VIIB:nde:136}证明与 \SGARefLink{sga3:VIIB:2.6.3:prop}{2.6.3} 的证明相同。}
可立即归约到 \(k\) 为 Artin 环的情形。置
\[
  L=\mathbf L(k).
\]
由 \SGARefLink{sga3:VIIB:2.3.1}{2.3.1}，
\[
  \operatorname{Hom}_{\operatorname{Grf}/k}
    \bigl(\mathscr G_p(\mathbf L),G\bigr)
\]
与使下图交换的含幺代数态射
\[
  h:U_p(L)\longrightarrow\mathbf H(G)
\]
的集合存在双射：
\begin{center}
  \phantomsection
  \label{sga3:VIIB:4.2.2:restricted-universal-diagram}
  % SGA3-VII-DIAG-122
  % SGA3-VII-DIAG-122
% Authority: Exp7B-13oct24.pdf, local p. 76.
% Crop: figures/Exp7B_localp076_restricted_enveloping_universal_diagram.png
% Locator: tex/components/68_expose_VIIB_section422_restricted_universal_property_en.tex, Loop 1 line 102.
\begin{tikzcd}[row sep=1.2em,column sep=3.9em]
U_p(L) \arrow[r,"h"] \arrow[dd,"\Delta_{U_p(L)}"']
  & \mathbf H(G) \arrow[dr,"\delta_G"] & {} \\
{} & {} & \mathbf H(G\times G) \\
U_p(L)\otimes U_p(L) \arrow[r,"h\otimes h"]
  & \mathbf H(G)\otimes\mathbf H(G)
      \arrow[ur,"\varphi_G"']
  & {.}
\end{tikzcd}

\end{center}
态射 \(h\) 由其在 \(L\) 上的限制决定，这个限制是从 \(L\) 到 \(\mathbf H(G)\) 的底 \(p\)-Lie 代数的 \(p\)-Lie 代数态射。该图交换，恰好是说 \(h\) 将 \(L\) 映入 \(\mathbf H(G)\) 中满足
\[
  \delta_G(x)
  =
  \varphi_G(x\otimes1+1\otimes x).
\]
的元素 \(x\) 构成的子集。由 \SGARefLink{sga3:VIIB:2.6}{2.6}\textup{(c)}，这个子集就是 \(\operatorname{Lie}(G)\)。命题得证。
